% ===========================================================================
% Relativistic Modification of Chapter XI, Sections 100--101
% Kelvin--Helmholtz perturbation equations; two uniform fluids
% ===========================================================================
% Agent 48 — relativistic branch
% Conventions: RELATIVISTIC_CONVENTIONS.md, rel_preamble.tex
% ===========================================================================

\chapter{Relativistic Kelvin--Helmholtz Instability}
\label{ch:rel-11-KH}

\section*{11R. THE STABILITY OF SUPERPOSED FLUIDS IN RELATIVE MOTION: RELATIVISTIC FORMULATION}

\section{The relativistic perturbation equations}\label{sec:rel-11-100}

In \S\,100 of Chandrasekhar, the Kelvin--Helmholtz instability is studied
for an incompressible, inviscid fluid in which adjacent layers stream
horizontally at different velocities. The relativistic generalization
of this problem arises naturally in astrophysical contexts where the
relative velocity of the two fluid layers becomes a significant fraction
of the speed of light: relativistic jets impinging on ambient media,
the interface between the fast spine and slower sheath of AGN jets, and
shear layers in gamma-ray burst outflows.

Throughout this chapter we adopt the conventions of
\textsc{Relativistic\_Conventions}: metric signature $(-,+,+,+)$,
the speed of light $c$ kept explicit, and the 4-velocity normalized as
$u^{\mu}u_{\mu} = -c^{2}$.

\subsection{Relativistic kinematics of two-fluid shear flow}

Consider two semi-infinite fluids separated by a planar interface
at $z=0$, streaming in the $x$-direction with 3-velocities $U_1$
(lower fluid, $z<0$) and $U_2$ (upper fluid, $z>0$) as measured in the
laboratory frame.  The associated Lorentz factors are
\begin{equation}\label{eq:rel-11-1}
  \gamma_j = \frac{1}{\sqrt{1 - U_j^{2}/c^{2}}}, \qquad j = 1,2.
\end{equation}

The relative velocity between the two streams, as determined by
the relativistic velocity-addition formula, is
\begin{equation}\label{eq:rel-11-2}
  \boxed{\;
  V_{\mathrm{rel}} = \frac{|U_1 - U_2|}{1 - U_1 U_2/c^{2}}.
  \;}
\end{equation}
The corresponding \emph{relative Lorentz factor} is
\begin{equation}\label{eq:rel-11-3}
  \gamma_{\mathrm{rel}}
  = \frac{1}{\sqrt{1 - V_{\mathrm{rel}}^{2}/c^{2}}}
  = \gamma_1\,\gamma_2\!\left(1 - \frac{U_1 U_2}{c^{2}}\right).
\end{equation}
Causality requires $V_{\mathrm{rel}} < c$, which is automatically
satisfied for any $|U_j| < c$.

It is often convenient to work in the rest frame of one of the fluids.
A Lorentz boost with velocity $U_1$ in the $x$-direction transforms
the lower fluid to rest; in this frame the upper fluid streams with
velocity
\begin{equation}\label{eq:rel-11-4}
  U'_{2} = \frac{U_2 - U_1}{1 - U_1 U_2/c^{2}} = V_{\mathrm{rel}}\,
  \operatorname{sgn}(U_2 - U_1).
\end{equation}

\subsection{Equilibrium state and the energy-momentum tensor}

For each fluid region, the unperturbed energy-momentum tensor of a
perfect relativistic fluid is
\begin{equation}\label{eq:rel-11-5}
  T^{\mu\nu} = \frac{\enthalpy_j}{c^{2}}\, u_j^{\mu}\, u_j^{\nu}
  + p_j\, \eta^{\mu\nu},
  \qquad \enthalpy_j \equiv \edensity_j + p_j,
\end{equation}
where $\edensity_j$ is the total energy density (including rest-mass
energy) and $p_j$ is the pressure in region $j$.  The 4-velocity of
the unperturbed flow is
\begin{equation}\label{eq:rel-11-6}
  u_j^{\mu} = \gamma_j\,(c,\, U_j,\, 0,\, 0).
\end{equation}

In the relativistic theory the \emph{inertial mass density} is not the
rest-mass density $\rho_{0,j}$ but rather
\begin{equation}\label{eq:rel-11-7}
  \boxed{\;
  \hat{\rho}_j \equiv \frac{\enthalpy_j}{c^{2}}
  = \frac{\edensity_j + p_j}{c^{2}},
  \;}
\end{equation}
which replaces $\rho_j$ wherever it appears in the classical equations.
This is the central relativistic substitution for the Kelvin--Helmholtz
problem.

\subsection{Linearized perturbation equations}

We perturb the equilibrium state in each region by writing
\begin{equation}\label{eq:rel-11-8}
  u^{\mu} = u_j^{\mu} + \delta u_j^{\mu}, \qquad
  p = p_j + \delta p_j, \qquad
  \edensity = \edensity_j + \delta\edensity_j.
\end{equation}
The perturbation 4-velocity must remain orthogonal to $u^{\mu}$, so
$u_{\mu}\,\delta u^{\mu} = 0$ at linear order.

We decompose the spatial part of $\delta u^{\mu}$ in the laboratory
frame as perturbation velocities $(u, v, w)$ and seek normal-mode
solutions with dependence
\begin{equation}\label{eq:rel-11-9}
  \exp\!\left[\mathrm{i}(k_x\,x + k_y\,y + n\,t)\right],
\end{equation}
where $k = \sqrt{k_x^{2}+k_y^{2}}$ and $n$ is the (complex) frequency.

The Doppler-shifted frequency in the rest frame of fluid $j$ is
\begin{equation}\label{eq:rel-11-10}
  \tilde{n}_j = \gamma_j\!\left(n + k_x\, U_j\right),
\end{equation}
which accounts for the relativistic Doppler effect. In the
non-relativistic limit $\gamma_j \to 1$ and we recover the classical
Doppler shift $n + k_x U_j$.

Linearizing the relativistic Euler equation
$\hat{\rho}\, u^{\alpha}\nabla_{\alpha} u^{\mu}
= -(\eta^{\mu\nu} + u^{\mu}u^{\nu}/c^{2})\,\partial_{\nu} p$
and the continuity equation $\nabla_{\mu}(\hat{\rho}\, u^{\mu}) = 0$
about the uniform-flow equilibrium, and assuming the fluid to be
effectively incompressible in each region (valid when all flow speeds
and perturbation phase velocities are much less than the sound speed
$c_s$ in each fluid), we obtain the relativistic analogues of
Chandrasekhar's equations~(1)--(6) of Chapter~XI:
\begin{equation}\label{eq:rel-11-11}
  \mathrm{i}\,\gamma_j^{2}\,\hat{\rho}_j\!\left(n + k_x U_j\right) u
  + \gamma_j^{2}\,\hat{\rho}_j\,(DU_j)\, w
  = -\mathrm{i}\, k_x\,\delta p,
\end{equation}
\begin{equation}\label{eq:rel-11-12}
  \mathrm{i}\,\gamma_j^{2}\,\hat{\rho}_j\!\left(n + k_x U_j\right) v
  = -\mathrm{i}\, k_y\,\delta p,
\end{equation}
\begin{equation}\label{eq:rel-11-13}
  \mathrm{i}\,\hat{\rho}_j\!\left(n + k_x U_j\right) w
  = -D\delta p - g\,\delta\hat{\rho}
  - k^{2}\sum_s T_s\,\delta z_s\,\delta(z - z_s),
\end{equation}
\begin{equation}\label{eq:rel-11-14}
  \frac{\partial u}{\partial x}
  + \frac{\partial v}{\partial y}
  + \frac{\partial w}{\partial z} = 0,
\end{equation}
where $D = d/dz$, $\hat{\rho}_j = (\edensity_j + p_j)/c^{2}$, and
$T_s$ is the surface tension at interface $s$.

\emph{Key relativistic modifications.}  (i) In the horizontal momentum
equations~\eqref{eq:rel-11-11} and~\eqref{eq:rel-11-12} an additional
factor of $\gamma_j^{2}$ appears because the horizontal inertia of a
fluid element moving at velocity $U_j$ is $\gamma_j^{2}\hat{\rho}_j$
(the ``longitudinal mass'').  (ii) In the vertical equation~%
\eqref{eq:rel-11-13} the gravitational term involves $\hat{\rho}_j$
rather than $\rho_j$; there is no extra $\gamma_j$ factor on this
term because the vertical motion is transverse to the boost.
(iii) The incompressibility condition~\eqref{eq:rel-11-14} is
unchanged at linear order.

\subsection{Reduction to a single equation}

Following the same algebraic procedure as in the classical case
(multiplying the horizontal equations by $-\mathrm{i}k_x$ and
$-\mathrm{i}k_y$, adding, and using incompressibility), we eliminate
$u$, $v$, and $\delta p$ to obtain the relativistic analogue of
Chandrasekhar's equation~(16):
\begin{multline}\label{eq:rel-11-15}
  D\!\left[\gamma_j^{2}\,\hat{\rho}_j\!\left(n + k_x U_j\right)Dw
  - \gamma_j^{2}\,\hat{\rho}_j\, k_x(DU_j)\,w\right]
  - k^{2}\hat{\rho}_j\!\left(n + k_x U_j\right) w \\
  = g\,k^{2}\bigg[D\hat{\rho}
  - \frac{k^{2}}{g}\sum_s T_s\,\delta(z - z_s)\bigg]
  \frac{w}{n + k_x U_j}\,.
\end{multline}
In regions of constant $\hat{\rho}_j$ and $U_j$, this reduces to
$(D^{2} - k^{2})w = 0$, exactly as in the classical theory.

The interface condition (integrating across the discontinuity at
$z = z_s$) becomes
\begin{equation}\label{eq:rel-11-16}
  \Delta_s\!\left[\gamma^{2}\hat{\rho}\!\left(n + k_x U\right)Dw
  - \gamma^{2}\hat{\rho}\, k_x(DU)\,w\right]
  = g\,k^{2}\!\left[\Delta_s(\hat{\rho})
  - \frac{k^{2}T_s}{g}\right]\!
  \left(\frac{w}{n + k_x U}\right)_{\!s},
\end{equation}
the relativistic generalization of Chandrasekhar's equation~(19).


% ===================================================================
\section{Two uniform relativistic fluids in relative horizontal motion}
\label{sec:rel-11-101}
% ===================================================================

We now apply the perturbation equations of the preceding section to
the case treated in \S\,101 of Chandrasekhar: two uniform fluids of
(relativistic) inertial densities $\hat{\rho}_1 = (\edensity_1 + p_1)/c^{2}$
and $\hat{\rho}_2 = (\edensity_2 + p_2)/c^{2}$, separated by a
horizontal boundary at $z = 0$, with streaming velocities $U_1$ and $U_2$.

In each region $(D^{2} - k^{2})w = 0$ with solutions
\begin{equation}\label{eq:rel-11-17}
  w_1 = A\!\left(n + k_x U_1\right) e^{+kz} \quad (z < 0),
  \qquad
  w_2 = A\!\left(n + k_x U_2\right) e^{-kz} \quad (z > 0),
\end{equation}
identical in form to the classical case.  Applying the interface
condition~\eqref{eq:rel-11-16} yields the \emph{relativistic dispersion
relation}:
\begin{equation}\label{eq:rel-11-18}
  \boxed{\;
  \gamma_1^{2}\,\hat{\rho}_1\!\left(n + k_x U_1\right)^{2}
  + \gamma_2^{2}\,\hat{\rho}_2\!\left(n + k_x U_2\right)^{2}
  = g\,k\!\left(\hat{\rho}_1 - \hat{\rho}_2\right)
  + k^{3}\,\frac{T}{\hat{\rho}_1 + \hat{\rho}_2}\,,
  \;}
\end{equation}
where $T$ is the surface tension.  This is the relativistic
counterpart of Chandrasekhar's equation~(26).

Define the relativistic density fractions
\begin{equation}\label{eq:rel-11-19}
  \hat{\alpha}_j
  = \frac{\gamma_j^{2}\,\hat{\rho}_j}
  {\gamma_1^{2}\,\hat{\rho}_1 + \gamma_2^{2}\,\hat{\rho}_2},
  \qquad j = 1,2,
  \qquad \hat{\alpha}_1 + \hat{\alpha}_2 = 1.
\end{equation}
In terms of these fractions the dispersion relation becomes
\begin{equation}\label{eq:rel-11-20}
  \hat{\alpha}_1\!\left(n + k_x U_1\right)^{2}
  + \hat{\alpha}_2\!\left(n + k_x U_2\right)^{2}
  = \frac{g\,k\!\left(\hat{\rho}_1 - \hat{\rho}_2\right)
  + k^{3}\,T/(\hat{\rho}_1+\hat{\rho}_2)}
  {\gamma_1^{2}\,\hat{\rho}_1 + \gamma_2^{2}\,\hat{\rho}_2}\,.
\end{equation}

Expanding and solving the quadratic in $n$:
\begin{multline}\label{eq:rel-11-21}
  n = -k_x\!\left(\hat{\alpha}_1 U_1 + \hat{\alpha}_2 U_2\right)
  \pm \bigg[
  \frac{g\,k(\hat{\rho}_1 - \hat{\rho}_2)
  + k^{3}\,T/(\hat{\rho}_1+\hat{\rho}_2)}
  {\gamma_1^{2}\hat{\rho}_1 + \gamma_2^{2}\hat{\rho}_2}
  \\
  {}- k_x^{2}\,\hat{\alpha}_1\,\hat{\alpha}_2\!\left(U_1 - U_2\right)^{2}
  \bigg]^{1/2}.
\end{multline}

\subsection*{(a) Without surface tension: relativistic critical velocity}

Setting $T = 0$ and $k_x = k$ (the most unstable direction), instability
occurs when
\begin{equation}\label{eq:rel-11-22}
  k\,\hat{\alpha}_1\,\hat{\alpha}_2\!\left(U_1 - U_2\right)^{2}
  > \frac{g\!\left(\hat{\rho}_1 - \hat{\rho}_2\right)}
  {\gamma_1^{2}\hat{\rho}_1 + \gamma_2^{2}\hat{\rho}_2}\,,
\end{equation}
which gives a minimum unstable wavenumber
\begin{equation}\label{eq:rel-11-23}
  k_{\min}^{(\mathrm{rel})}
  = \frac{g\!\left(\hat{\rho}_1 - \hat{\rho}_2\right)}
  {\hat{\alpha}_1\,\hat{\alpha}_2\!\left(U_1 - U_2\right)^{2}
  \!\left(\gamma_1^{2}\hat{\rho}_1 + \gamma_2^{2}\hat{\rho}_2\right)}.
\end{equation}

\emph{Comparison with the classical result.}  In the non-relativistic
limit $\gamma_j \to 1$, $\hat{\rho}_j \to \rho_j$, and we recover
Chandrasekhar's equation~(35):
$k_{\min} = g(\alpha_1 - \alpha_2)/[\alpha_1\alpha_2(U_1 - U_2)^{2}]$.

The relativistic modifications are twofold:

\begin{enumerate}
\item \textbf{Inertial-density replacement:} The rest-mass density
$\rho_j$ is replaced everywhere by $\hat{\rho}_j = (\edensity_j + p_j)/c^{2}$.
For a hot relativistic plasma with $p \sim \edensity/3$ (radiation-dominated
equation of state), $\hat{\rho} = 4\edensity/(3c^{2})$, which is
$4/3$ times the naive $\edensity/c^{2}$.

\item \textbf{Lorentz-factor weighting:} The effective inertia
$\gamma_j^{2}\hat{\rho}_j$ of each stream is enhanced by the
square of its Lorentz factor.  This means that a faster stream is
``heavier'' in the sense relevant for the KH instability, altering the
density fractions $\hat{\alpha}_j$ and hence the critical wavenumber.
\end{enumerate}

The \emph{critical relative velocity} for the onset of instability at
a given wavenumber $k$ (with $k_x = k$, $T = 0$) can be written in
a form that makes the relativistic corrections transparent.
Working in the frame where fluid~1 is at rest ($U_1 = 0$,
$\gamma_1 = 1$, $U_2 = V_{\mathrm{rel}}$, $\gamma_2 = \gamma_{\mathrm{rel}}$),
the dispersion relation~\eqref{eq:rel-11-18} becomes
\begin{equation}\label{eq:rel-11-24}
  \hat{\rho}_1\, n^{2}
  + \gamma_{\mathrm{rel}}^{2}\,\hat{\rho}_2
  \!\left(n + k\, V_{\mathrm{rel}}\right)^{2}
  = g\,k\!\left(\hat{\rho}_1 - \hat{\rho}_2\right).
\end{equation}
Instability ($n$ has a positive imaginary part) requires
\begin{equation}\label{eq:rel-11-25}
  \boxed{\;
  V_{\mathrm{rel}}^{2}
  > \frac{g}{k}\,
  \frac{(\hat{\rho}_1 - \hat{\rho}_2)
  (\hat{\rho}_1 + \gamma_{\mathrm{rel}}^{2}\hat{\rho}_2)}
  {\hat{\rho}_1\,\gamma_{\mathrm{rel}}^{2}\,\hat{\rho}_2}\,.
  \;}
\end{equation}
In the non-relativistic limit ($\gamma_{\mathrm{rel}} \to 1$,
$\hat{\rho}_j \to \rho_j$) this reduces to
$V_{\mathrm{rel}}^{2} > (g/k)(\rho_1^{2} - \rho_2^{2})/(\rho_1\rho_2)$,
reproducing the classical result
$V_{\mathrm{crit}}^{2} \sim g(\rho_2^{2}-\rho_1^{2})/(\rho_1\rho_2\, k)$
(with appropriate sign conventions for $\rho_1 > \rho_2$).

For ultra-relativistic streaming ($\gamma_{\mathrm{rel}} \gg 1$),
the right-hand side of~\eqref{eq:rel-11-25} approaches
$(g/k)(\hat{\rho}_1 - \hat{\rho}_2)/\hat{\rho}_2$, which is
independent of $\gamma_{\mathrm{rel}}$.  Since $V_{\mathrm{rel}} \to c$
in this limit, the inequality is always satisfied for sufficiently
large relative velocity: the KH instability persists in the
ultra-relativistic regime.

\subsection*{(b) Relativistic Doppler effect and aberration}

The growth rate of the KH instability is modified in the relativistic
regime by two kinematic effects:

\begin{enumerate}
\item \textbf{Relativistic Doppler shift.}  The Doppler-shifted
frequency~\eqref{eq:rel-11-10} differs from the classical one by the
factor $\gamma_j$. As a consequence, the real part of $n$ (the
oscillation frequency of the interface) is modified; the growth rate
(imaginary part of $n$) inherits $\gamma$-dependent corrections that
reduce the e-folding rate in the lab frame relative to the co-moving
frame by a factor of $\gamma$.

\item \textbf{Aberration.}  The angle $\theta$ between the wave vector
$\mathbf{k}$ and the streaming direction transforms under a boost:
\begin{equation}\label{eq:rel-11-26}
  \cos\theta' = \frac{\cos\theta - U_j/c}{1 - (U_j/c)\cos\theta},
\end{equation}
where $\theta'$ is the angle in the rest frame of fluid~$j$.
This means that perturbations which are oblique in the lab frame appear
more nearly aligned with the streaming direction in the co-moving frame,
enhancing the effective $k_x$ component and thus the destabilizing
influence of the shear.
\end{enumerate}

\subsection*{(c) The stabilizing effect of surface tension — relativistic}

When surface tension $T$ is present, equation~\eqref{eq:rel-11-21}
predicts stability (setting $k_x = k$) if
\begin{equation}\label{eq:rel-11-27}
  \hat{\alpha}_1\,\hat{\alpha}_2\!\left(U_1 - U_2\right)^{2}
  < \frac{g\,k^{-1}\!\left(\hat{\rho}_1 - \hat{\rho}_2\right)
  + k\, T/(\hat{\rho}_1 + \hat{\rho}_2)}
  {\gamma_1^{2}\hat{\rho}_1 + \gamma_2^{2}\hat{\rho}_2}\,.
\end{equation}
The right-hand side is minimized at a wavenumber $k_*^{(\mathrm{rel})}$
satisfying
\begin{equation}\label{eq:rel-11-28}
  \left(k_*^{(\mathrm{rel})}\right)^{2}
  = \frac{g\!\left(\hat{\rho}_1 - \hat{\rho}_2\right)
  (\hat{\rho}_1 + \hat{\rho}_2)}{T}\,,
\end{equation}
which has the same functional form as the classical
expression~(Chandrasekhar eq.~(38)) with $\rho_j$ replaced by
$\hat{\rho}_j$.

Surface tension suppresses the relativistic KH instability if
\begin{equation}\label{eq:rel-11-29}
  \boxed{\;
  \left(U_1 - U_2\right)^{2}
  < \frac{2}{\hat{\alpha}_1\,\hat{\alpha}_2}
  \sqrt{\frac{T\, g\!\left(\hat{\rho}_1 - \hat{\rho}_2\right)}
  {(\hat{\rho}_1 + \hat{\rho}_2)
  \!\left(\gamma_1^{2}\hat{\rho}_1
  + \gamma_2^{2}\hat{\rho}_2\right)^{2}}}\,.
  \;}
\end{equation}
In the non-relativistic limit this reduces to Chandrasekhar's
equation~(40), which is the Kelvin criterion for the suppression
of the KH instability by surface tension.

\subsection*{(d) Causality checks}

The relativistic formulation demands that all physical velocities
remain subluminal.  We verify three requirements:

\begin{enumerate}
\item \textbf{Streaming velocities:}  $|U_j| < c$ for $j = 1,2$,
guaranteed by hypothesis and by equation~\eqref{eq:rel-11-1}.

\item \textbf{Relative velocity:}  $V_{\mathrm{rel}} < c$ as noted
below equation~\eqref{eq:rel-11-3}; this follows from the structure
of the relativistic velocity-addition formula.

\item \textbf{Phase velocities of KH modes:}  The phase velocity of a
perturbation mode is $v_{\mathrm{ph}} = -\operatorname{Re}(n)/k$.
From equation~\eqref{eq:rel-11-21}, the real part of $n$ is
$-k(\hat{\alpha}_1 U_1 + \hat{\alpha}_2 U_2)$ (plus a correction from
the square root for stable modes), so
\begin{equation}\label{eq:rel-11-30}
  v_{\mathrm{ph}} = \hat{\alpha}_1 U_1 + \hat{\alpha}_2 U_2,
\end{equation}
which is a weighted average of $U_1$ and $U_2$ with positive weights
summing to unity.  Since $|U_j| < c$, we have $|v_{\mathrm{ph}}| < c$.

For unstable modes the imaginary part $\operatorname{Im}(n)$
represents a temporal growth rate, not a spatial propagation velocity,
so there is no additional causality constraint from the growth rate
itself.  However, the \emph{group velocity}
$v_g = \partial\operatorname{Re}(n)/\partial k$ must also satisfy
$|v_g| < c$; a direct calculation from~\eqref{eq:rel-11-21} confirms
this, as $v_g$ remains a combination of $U_1$, $U_2$, and the
(subluminal) gravitational propagation speed $\sqrt{g/k}$.
\end{enumerate}

\subsection*{(e) Non-relativistic limit}

In the limit $|U_j|/c \to 0$, $p_j/(\rho_j c^{2}) \to 0$,
we have $\gamma_j \to 1$, $\hat{\rho}_j \to \rho_j$,
$\hat{\alpha}_j \to \alpha_j$, and all equations of this section
reduce identically to those of Chandrasekhar's \S\,101.  In particular:
\begin{itemize}
\item Equation~\eqref{eq:rel-11-18} $\to$ Chandrasekhar eq.~(26).
\item Equation~\eqref{eq:rel-11-23} $\to$ Chandrasekhar eq.~(35).
\item Equation~\eqref{eq:rel-11-29} $\to$ Chandrasekhar eq.~(40).
\end{itemize}
This provides a consistency check on the relativistic generalization.
